\chapter{Overview of potentially helpful software}
\chaptershorttitle{Helpful software}
\label{cha:software}

\section{Overview of open-source programs and tools}
\begin{itemize}
	\item \href{https://github.com/ElmerCSC/elmerfem}{Elmer FEM} (Linux, macOS, Windows): FEM simulation tool
	\item \href{https://github.com/upb-lea/FEM_Magnetics_Toolbox}{FEM Magnetics Toolbox}: automated inductor and transformer design tool
	\item \href{https://en.wikipedia.org/wiki/FreeCAD}{FreeCAD} (Linux, macOS, Windows): 3D CAD software
	\item \href{https://github.com/geckocircuits/GeckoCIRCUITS}{GeckoCIRCUITS} (Linux, macOS, Windows): circuit simulator optimized for power electronics
	\item \href{https://en.wikipedia.org/wiki/GIMP}{GIMP} (Linux, macOS, Windows): raster graphics editor
	\item \href{https://git-scm.com/}{Git} (Linux, macOS, Windows): distributed version control system\footnote{Keeping a version-controlled copy of your thesis manuscript on an external system is highly recommended to protect against data loss if local hardware fails.}
	\item \href{https://en.wikipedia.org/wiki/Inkscape}{Inkscape} (Linux, macOS, Windows): vector graphics editor
	\item \href{https://github.com/upb-lea/Inkscape_electric_Symbols}{Inkscape Electrical Symbols}: electrical symbols library for Inkscape
	\item \href{https://en.wikipedia.org/wiki/JabRef}{JabRef} (Linux, macOS, Windows): reference management software
	\item \href{https://julialang.org/}{Julia} (Linux, macOS, Windows): programming language for high-performance numerical and scientific computing
	\item \href{https://jupyter.org/}{JupyterLab} (Linux, macOS, Windows): interactive environment for computational notebooks, code, data, and visualizations
	\item \href{https://en.wikipedia.org/wiki/KiCad}{KiCad} (Linux, macOS, Windows): electronic design and simulation tool for PCBs
	\item \href{https://en.wikipedia.org/wiki/GNU_Octave}{Octave} (Linux, macOS, Windows): MATLAB-like numerical computing environment
	\item \href{https://onelab.info/}{ONELAB} (Linux, macOS, Windows): FEM simulation tool
	\item \href{https://openfoam.org/}{OpenFOAM} (Linux; macOS and Windows via virtualized Linux environments): CFD simulation toolkit
	\item \href{https://github.com/OpenModelica/OpenModelica}{OpenModelica} (Linux, macOS, Windows): component-oriented systems modeling environment
	\item \href{https://www.python.org/}{Python} (Linux, macOS, Windows): general-purpose programming language with a broad scientific-computing ecosystem
	\item \href{https://en.wikipedia.org/wiki/TeXstudio}{TeXstudio} (Linux, macOS, Windows): \LaTeX\ editor
	\item \href{https://github.com/upb-lea/transistordatabase}{Transistordatabase}: transistor data management tool
\end{itemize}

\section{Other programs with proprietary licenses}
\begin{itemize}
	\item \href{https://en.wikipedia.org/wiki/Altium_Designer}{Altium Designer} (Windows): PCB and electronic design automation software
	\item \href{https://www.ansys.com/de-de/academic/free-student-products}{Ansys Student} (Windows): large-scale design and simulation toolkit
	\item \href{https://en.wikipedia.org/wiki/MATLAB}{MATLAB/Simulink} (Linux, macOS, Windows): general-purpose scientific computing
	\item \href{https://en.wikipedia.org/wiki/PyCharm}{PyCharm} (Linux, macOS, Windows): Python IDE
	\item \href{https://en.wikipedia.org/wiki/SolidWorks}{SolidWorks} (Windows): 3D CAD software
	\item \href{https://en.wikipedia.org/wiki/Visual_Studio_Code}{Visual Studio Code} (Linux, macOS, Windows): general-purpose editor (incl. \LaTeX\ support)\footnote{A freely licensed alternative is \href{https://vscodium.com/}{VSCodium}, a community-built distribution of the open-source VS Code codebase.}
	\item \href{https://en.wikipedia.org/wiki/Wolfram_Mathematica}{Wolfram Mathematica} (Linux, macOS, Windows): general-purpose scientific computing
\end{itemize}
%
